\IfFileExists{stacks-project.cls}{%
\documentclass{stacks-project}
}{%
\documentclass{amsart}
}

% For dealing with references we use the comment environment
\usepackage{verbatim}
\newenvironment{reference}{\comment}{\endcomment}
%\newenvironment{reference}{}{}
\newenvironment{slogan}{\comment}{\endcomment}
\newenvironment{history}{\comment}{\endcomment}

% For commutative diagrams we use Xy-pic
\usepackage[all]{xy}

% We use 2cell for 2-commutative diagrams.
\xyoption{2cell}
\UseAllTwocells

% We use multicol for the list of chapters between chapters
\usepackage{multicol}

% This is generally recommended for better output
\usepackage{lmodern}
\usepackage[T1]{fontenc}

% For cross-file-references
\usepackage{xr-hyper}

% Package for hypertext links:
\usepackage{hyperref}

% For any local file, say "hello.tex" you want to link to please
% use \externaldocument[hello-]{hello}
\externaldocument[introduction-]{introduction}
\externaldocument[conventions-]{conventions}
\externaldocument[sets-]{sets}
\externaldocument[categories-]{categories}
\externaldocument[topology-]{topology}
\externaldocument[sheaves-]{sheaves}
\externaldocument[sites-]{sites}
\externaldocument[stacks-]{stacks}
\externaldocument[fields-]{fields}
\externaldocument[algebra-]{algebra}
\externaldocument[brauer-]{brauer}
\externaldocument[homology-]{homology}
\externaldocument[derived-]{derived}
\externaldocument[simplicial-]{simplicial}
\externaldocument[more-algebra-]{more-algebra}
\externaldocument[smoothing-]{smoothing}
\externaldocument[modules-]{modules}
\externaldocument[sites-modules-]{sites-modules}
\externaldocument[injectives-]{injectives}
\externaldocument[cohomology-]{cohomology}
\externaldocument[sites-cohomology-]{sites-cohomology}
\externaldocument[dga-]{dga}
\externaldocument[dpa-]{dpa}
\externaldocument[sdga-]{sdga}
\externaldocument[hypercovering-]{hypercovering}
\externaldocument[schemes-]{schemes}
\externaldocument[constructions-]{constructions}
\externaldocument[properties-]{properties}
\externaldocument[morphisms-]{morphisms}
\externaldocument[coherent-]{coherent}
\externaldocument[divisors-]{divisors}
\externaldocument[limits-]{limits}
\externaldocument[varieties-]{varieties}
\externaldocument[topologies-]{topologies}
\externaldocument[descent-]{descent}
\externaldocument[perfect-]{perfect}
\externaldocument[more-morphisms-]{more-morphisms}
\externaldocument[flat-]{flat}
\externaldocument[groupoids-]{groupoids}
\externaldocument[more-groupoids-]{more-groupoids}
\externaldocument[etale-]{etale}
\externaldocument[chow-]{chow}
\externaldocument[intersection-]{intersection}
\externaldocument[pic-]{pic}
\externaldocument[weil-]{weil}
\externaldocument[adequate-]{adequate}
\externaldocument[dualizing-]{dualizing}
\externaldocument[duality-]{duality}
\externaldocument[discriminant-]{discriminant}
\externaldocument[derham-]{derham}
\externaldocument[local-cohomology-]{local-cohomology}
\externaldocument[algebraization-]{algebraization}
\externaldocument[curves-]{curves}
\externaldocument[resolve-]{resolve}
\externaldocument[models-]{models}
\externaldocument[functors-]{functors}
\externaldocument[equiv-]{equiv}
\externaldocument[pione-]{pione}
\externaldocument[etale-cohomology-]{etale-cohomology}
\externaldocument[proetale-]{proetale}
\externaldocument[relative-cycles-]{relative-cycles}
\externaldocument[more-etale-]{more-etale}
\externaldocument[trace-]{trace}
\externaldocument[crystalline-]{crystalline}
\externaldocument[spaces-]{spaces}
\externaldocument[spaces-properties-]{spaces-properties}
\externaldocument[spaces-morphisms-]{spaces-morphisms}
\externaldocument[decent-spaces-]{decent-spaces}
\externaldocument[spaces-cohomology-]{spaces-cohomology}
\externaldocument[spaces-limits-]{spaces-limits}
\externaldocument[spaces-divisors-]{spaces-divisors}
\externaldocument[spaces-over-fields-]{spaces-over-fields}
\externaldocument[spaces-topologies-]{spaces-topologies}
\externaldocument[spaces-descent-]{spaces-descent}
\externaldocument[spaces-perfect-]{spaces-perfect}
\externaldocument[spaces-more-morphisms-]{spaces-more-morphisms}
\externaldocument[spaces-flat-]{spaces-flat}
\externaldocument[spaces-groupoids-]{spaces-groupoids}
\externaldocument[spaces-more-groupoids-]{spaces-more-groupoids}
\externaldocument[bootstrap-]{bootstrap}
\externaldocument[spaces-pushouts-]{spaces-pushouts}
\externaldocument[spaces-chow-]{spaces-chow}
\externaldocument[groupoids-quotients-]{groupoids-quotients}
\externaldocument[spaces-more-cohomology-]{spaces-more-cohomology}
\externaldocument[spaces-simplicial-]{spaces-simplicial}
\externaldocument[spaces-duality-]{spaces-duality}
\externaldocument[formal-spaces-]{formal-spaces}
\externaldocument[restricted-]{restricted}
\externaldocument[spaces-resolve-]{spaces-resolve}
\externaldocument[formal-defos-]{formal-defos}
\externaldocument[defos-]{defos}
\externaldocument[cotangent-]{cotangent}
\externaldocument[examples-defos-]{examples-defos}
\externaldocument[algebraic-]{algebraic}
\externaldocument[examples-stacks-]{examples-stacks}
\externaldocument[stacks-sheaves-]{stacks-sheaves}
\externaldocument[criteria-]{criteria}
\externaldocument[artin-]{artin}
\externaldocument[quot-]{quot}
\externaldocument[stacks-properties-]{stacks-properties}
\externaldocument[stacks-morphisms-]{stacks-morphisms}
\externaldocument[stacks-limits-]{stacks-limits}
\externaldocument[stacks-cohomology-]{stacks-cohomology}
\externaldocument[stacks-perfect-]{stacks-perfect}
\externaldocument[stacks-introduction-]{stacks-introduction}
\externaldocument[stacks-more-morphisms-]{stacks-more-morphisms}
\externaldocument[stacks-geometry-]{stacks-geometry}
\externaldocument[moduli-]{moduli}
\externaldocument[moduli-curves-]{moduli-curves}
\externaldocument[examples-]{examples}
\externaldocument[exercises-]{exercises}
\externaldocument[guide-]{guide}
\externaldocument[desirables-]{desirables}
\externaldocument[coding-]{coding}
\externaldocument[obsolete-]{obsolete}
\externaldocument[fdl-]{fdl}
\externaldocument[index-]{index}

% Theorem environments.
%
\theoremstyle{plain}
\newtheorem{theorem}[subsection]{Theorem}
\newtheorem{proposition}[subsection]{Proposition}
\newtheorem{lemma}[subsection]{Lemma}

\theoremstyle{definition}
\newtheorem{definition}[subsection]{Definition}
\newtheorem{example}[subsection]{Example}
\newtheorem{exercise}[subsection]{Exercise}
\newtheorem{situation}[subsection]{Situation}

\theoremstyle{remark}
\newtheorem{remark}[subsection]{Remark}
\newtheorem{remarks}[subsection]{Remarks}

\numberwithin{equation}{subsection}

% Macros
%
\def\lim{\mathop{\mathrm{lim}}\nolimits}
\def\colim{\mathop{\mathrm{colim}}\nolimits}
\def\Spec{\mathop{\mathrm{Spec}}}
\def\Hom{\mathop{\mathrm{Hom}}\nolimits}
\def\Ext{\mathop{\mathrm{Ext}}\nolimits}
\def\SheafHom{\mathop{\mathcal{H}\!\mathit{om}}\nolimits}
\def\SheafExt{\mathop{\mathcal{E}\!\mathit{xt}}\nolimits}
\def\Sch{\mathit{Sch}}
\def\Mor{\mathop{\mathrm{Mor}}\nolimits}
\def\Ob{\mathop{\mathrm{Ob}}\nolimits}
\def\Sh{\mathop{\mathit{Sh}}\nolimits}
\def\NL{\mathop{N\!L}\nolimits}
\def\CH{\mathop{\mathrm{CH}}\nolimits}
\def\proetale{{pro\text{-}\acute{e}tale}}
\def\etale{{\acute{e}tale}}
\def\QCoh{\mathit{QCoh}}
\def\Ker{\mathop{\mathrm{Ker}}}
\def\Im{\mathop{\mathrm{Im}}}
\def\Coker{\mathop{\mathrm{Coker}}}
\def\Coim{\mathop{\mathrm{Coim}}}

% Boxtimes
%
\DeclareMathSymbol{\boxtimes}{\mathbin}{AMSa}{"02}

%
% Macros for moduli stacks/spaces
%
\def\QCohstack{\mathcal{QC}\!\mathit{oh}}
\def\Cohstack{\mathcal{C}\!\mathit{oh}}
\def\Spacesstack{\mathcal{S}\!\mathit{paces}}
\def\Quotfunctor{\mathrm{Quot}}
\def\Hilbfunctor{\mathrm{Hilb}}
\def\Curvesstack{\mathcal{C}\!\mathit{urves}}
\def\Polarizedstack{\mathcal{P}\!\mathit{olarized}}
\def\Complexesstack{\mathcal{C}\!\mathit{omplexes}}
% \Pic is the operator that assigns to X its picard group, usage \Pic(X)
% \Picardstack_{X/B} denotes the Picard stack of X over B
% \Picardfunctor_{X/B} denotes the Picard functor of X over B
\def\Pic{\mathop{\mathrm{Pic}}\nolimits}
\def\Picardstack{\mathcal{P}\!\mathit{ic}}
\def\Picardfunctor{\mathrm{Pic}}
\def\Deformationcategory{\mathcal{D}\!\mathit{ef}}
\usepackage{cleveref}
\crefname{originaldefinition}{original definition}{original definitions}
\Crefname{originaldefinition}{Original definition}{Original definitions}
\makeatletter
\newlabel{definition-separated}{{0AJ7}{source}{}{sourcepointer-1}{}}
\newlabel{definition-separated@cref}{{[originaldefinition][0][]0AJ7}{[1][0][]source}{}{}{}}
\makeatother
\crefname{originallemma}{original lemma}{original lemmas}
\Crefname{originallemma}{Original lemma}{Original lemmas}
\makeatletter
\newlabel{lemma-base-change-separated}{{0ARP}{source}{}{sourcepointer-2}{}}
\newlabel{lemma-base-change-separated@cref}{{[originallemma][0][]0ARP}{[1][0][]source}{}{}{}}
\makeatother
\crefname{originallemma}{original lemma}{original lemmas}
\Crefname{originallemma}{Original lemma}{Original lemmas}
\makeatletter
\newlabel{lemma-factor-through-thickening}{{0AIA}{source}{}{sourcepointer-3}{}}
\newlabel{lemma-factor-through-thickening@cref}{{[originallemma][0][]0AIA}{[1][0][]source}{}{}{}}
\makeatother
\crefname{originallemma}{original lemma}{original lemmas}
\Crefname{originallemma}{Original lemma}{Original lemmas}
\makeatletter
\newlabel{spaces-lemma-descent-representable-transformations-property}{{03KD}{source}{}{sourcepointer-4}{}}
\newlabel{spaces-lemma-descent-representable-transformations-property@cref}{{[originallemma][0][]03KD}{[1][0][]source}{}{}{}}
\makeatother
\crefname{originalremark}{original remark}{original remarks}
\Crefname{originalremark}{Original remark}{Original remarks}
\makeatletter
\newlabel{remark-set-theoretic}{{0AIS}{source}{}{sourcepointer-5}{}}
\newlabel{remark-set-theoretic@cref}{{[originalremark][0][]0AIS}{[1][0][]source}{}{}{}}
\makeatother
\title{The separation axioms for a morphism of formal algebraic spaces can be checked on affine base changes}
\author{The Stacks Project Authors}
\date{stacks-project:formal-spaces.tex:local-snapshot-20261012; GNU FDL 1.2 or later}
\begin{document}
\maketitle
\noindent\textit{Editorial note.} All delivered mathematical text is the author's own, copied from the pinned Stacks Project chapter file formal-spaces.tex; no statement or proof step was rewritten, reordered or removed.\par\medskip
\noindent\textit{Editorial note.} The proof displays five commutative diagrams and formulas of fibre products; all lie inside the counted span and are reproduced verbatim.\par\medskip
\noindent\textit{Editorial note.} The counted spans carry 7 navigation stubs; no bibliography citation occurs in the spans.\par\medskip
\noindent\textit{Editorial note.} The preamble span points at the chapter's own preamble.tex shipped with the frozen source, used verbatim as staged.\par\medskip
\noindent\textit{Allowed standard prerequisites (editorial).}\begin{itemize}
\item Stability of the separation axioms under base change. Same chapter, Lemma lemma-base-change-separated, invoked at proof lines 6913-6915. Gives the two immediate implications.
\item The definition of the separation axioms for a morphism and of being of the corresponding type as a formal algebraic space. Same chapter, Definition definition-separated-morphism at statement lines 6887-6888, and Definition definition-separated at statement lines 6899-6900. Fix the properties whose equivalence is asserted.
\item Affine formal algebraic spaces as filtered colimits, with factorisation of morphisms into the colimit through a member. Same chapter, Definition definition-affine-formal-algebraic-space at proof lines 6956-6958 and Lemma lemma-factor-through-thickening at proof lines 6970-6972. Underlies the passage from the third to the fifth formulation.
\item The definition of a formal algebraic space through a presentation, and descent of representable transformation properties. Same chapter, Definition definition-formal-algebraic-space at proof lines 6917-6918 and 6974-6976, and Spaces, Lemma spaces-lemma-descent-representable-transformations-property at proof lines 6976-6981. Supply the presentation and transfer the property along the covering.
\end{itemize}\medskip
\ifdefined\hypertarget\hypertarget{sourcepointer-1}{}\fi
\noindent\textit{Source reference.} Original definition 0AJ7: sources/stacks-algebraic-stacks-morphisms/pinned-original-source.tex, source lines 601--601. Pointer outside the counted span, invoked once for context; the load-bearing argument of the counted statement is the author's proof reproduced in full in the proof span.\par\medskip
\ifdefined\hypertarget\hypertarget{sourcepointer-2}{}\fi
\noindent\textit{Source reference.} Original lemma 0ARP: sources/stacks-algebraic-stacks-morphisms/pinned-original-source.tex, source lines 703--703. Pointer outside the counted span, invoked once for context; the load-bearing argument of the counted statement is the author's proof reproduced in full in the proof span.\par\medskip
\ifdefined\hypertarget\hypertarget{sourcepointer-3}{}\fi
\noindent\textit{Source reference.} Original lemma 0AIA: sources/stacks-formal-spaces/pinned-original-source.tex, source lines 2080--2080. Pointer outside the counted span, invoked once for context; the load-bearing argument of the counted statement is the author's proof reproduced in full in the proof span.\par\medskip
\ifdefined\hypertarget\hypertarget{sourcepointer-4}{}\fi
\noindent\textit{Source reference.} Original lemma 03KD: sources/stacks-chow/tags.txt, source lines 4673--4673. Pointer outside the counted span, invoked once for context; the load-bearing argument of the counted statement is the author's proof reproduced in full in the proof span.\par\medskip
\ifdefined\hypertarget\hypertarget{sourcepointer-5}{}\fi
\noindent\textit{Source reference.} Original remark 0AIS: sources/stacks-formal-spaces/pinned-original-source.tex, source lines 2769--2769. Pointer outside the counted span, invoked once for context; the load-bearing argument of the counted statement is the author's proof reproduced in full in the proof span.\par\medskip
\begin{definition}
\label{definition-affine-formal-algebraic-space}
Let $S$ be a scheme. We say a sheaf $X$ on $(\Sch/S)_{fppf}$ is an
{\it affine formal algebraic space} if there exist
\begin{enumerate}
\item a directed set $\Lambda$,
\item a system $(X_\lambda, f_{\lambda \mu})$ over $\Lambda$
in $(\Sch/S)_{fppf}$ where
\begin{enumerate}
\item each $X_\lambda$ is affine,
\item each $f_{\lambda \mu} : X_\lambda \to X_\mu$ is a thickening,
\end{enumerate}
\end{enumerate}
such that
$$
X \cong \colim_{\lambda \in \Lambda} X_\lambda
$$
as fppf sheaves and $X$ satisfies a set theoretic condition
(see Remark \ref{remark-set-theoretic}). A
{\it morphism of affine formal algebraic spaces}
over $S$ is a map of sheaves.
\end{definition}

\begin{definition}
\label{definition-formal-algebraic-space}
Let $S$ be a scheme. We say a sheaf $X$ on $(\Sch/S)_{fppf}$ is a
{\it formal algebraic space} if there exist a family of maps
$\{X_i \to X\}_{i \in I}$ of sheaves such that
\begin{enumerate}
\item $X_i$ is an affine formal algebraic space,
\item $X_i \to X$ is representable by algebraic spaces and \'etale,
\item $\coprod X_i \to X$ is surjective as a map of sheaves
\end{enumerate}
and $X$ satisfies a set theoretic condition
(see Remark \ref{remark-set-theoretic}). A
{\it morphism of formal algebraic spaces}
over $S$ is a map of sheaves.
\end{definition}

\begin{definition}
\label{definition-separated-morphism}
Let $S$ be a scheme.
Let $f : X \to Y$ be a morphism of formal algebraic spaces over $S$.
Let $\Delta_{X/Y} : X \to X \times_Y X$ be the diagonal morphism.
\begin{enumerate}
\item We say $f$ is {\it separated} if $\Delta_{X/Y}$ is a closed immersion.
\item We say $f$ is {\it quasi-separated} if $\Delta_{X/Y}$ is quasi-compact.
\end{enumerate}
\end{definition}

\begin{lemma}
\label{lemma-separated-local}
Let $S$ be a scheme. Let $f : X \to Y$ be a morphism of formal algebraic spaces
over $S$. Let $\mathcal{P}$ be any of the separation axioms of
Definition \ref{definition-separated-morphism}.
The following are equivalent
\begin{enumerate}
\item $f$ is $\mathcal{P}$,
\item for every scheme $Z$ and morphism $Z \to Y$ the
base change $Z \times_Y X \to Z$ of $f$ is $\mathcal{P}$,
\item for every affine scheme $Z$ and every morphism $Z \to Y$ the
base change $Z \times_Y X \to Z$ of $f$ is $\mathcal{P}$,
\item for every affine scheme $Z$ and every morphism $Z \to Y$ the
formal algebraic space $Z \times_Y X$ is $\mathcal{P}$ (see
Definition \ref{definition-separated}),
\item there exists a covering $\{Y_j \to Y\}$ as in
Definition \ref{definition-formal-algebraic-space}
such that the base change $Y_j \times_Y X \to Y_j$ has
$\mathcal{P}$ for all $j$.
\end{enumerate}
\end{lemma}

\begin{proof}
We will repeatedly use
Lemma \ref{lemma-base-change-separated}
without further mention. In particular, it is clear that
(1) implies (2) and (2) implies (3).

\medskip\noindent
Assume (3) and let $Z \to Y$ be a morphism where $Z$ is an affine scheme.
Let $U$, $V$ be affine schemes and let $a : U \to Z \times_Y X$
and $b : V \to Z \times_Y X$ be morphisms. Then
$$
U \times_{Z \times_Y X} V =
(Z \times_Y X) \times_{\Delta, (Z \times_Y X) \times_Z (Z \times_Y X)}
(U \times_Z V)
$$
and we see that this is quasi-compact if $\mathcal{P} =$``quasi-separated''
or an affine scheme equipped with a closed immersion into
$U \times_Z V$ if $\mathcal{P} =$``separated''. Thus (4) holds.

\medskip\noindent
Assume (4) and let $Z \to Y$ be a morphism where $Z$ is an affine scheme.
Let $U$, $V$ be affine schemes and let $a : U \to Z \times_Y X$
and $b : V \to Z \times_Y X$ be morphisms. Reading the argument above
backwards, we see that $U \times_{Z \times_Y X} V \to U \times_Z V$
is quasi-compact if $\mathcal{P} =$``quasi-separated'' or a closed
immersion if $\mathcal{P} =$``separated''. Since we can choose $U$ and
$V$ as above such that $U$ varies through an
\'etale covering of $Z \times_Y X$, we find
that the corresponding morphisms
$$
U \times_Z V \to (Z \times_Y X) \times_Z (Z \times_Y X)
$$
form an \'etale covering by affines. Hence we conclude that
$\Delta : (Z \times_Y X) \to (Z \times_Y X) \times_Z (Z \times_Y X)$
is quasi-compact, resp.\ a closed immersion. Thus (3) holds.

\medskip\noindent
Let us prove that (3) implies (5). Assume (3) and let
$\{Y_j \to Y\}$ be as in
Definition \ref{definition-formal-algebraic-space}.
We have to show that the morphisms
$$
\Delta_j :
Y_j \times_Y X
\longrightarrow
(Y_j \times_Y X) \times_{Y_j} (Y_j \times_Y X) =
Y_j \times_Y X \times_Y X
$$
has the corresponding property (i.e., is quasi-compact or a closed immersion).
Write $Y_j = \colim Y_{j, \lambda}$ as in
Definition \ref{definition-affine-formal-algebraic-space}.
Replacing $Y_j$ by $Y_{j, \lambda}$ in the formula above, we have the
property by our assumption that (3) holds. Since the displayed arrow
is the colimit of the arrows $\Delta_{j, \lambda}$ and since we
can test whether $\Delta_j$ has the corresponding property by
testing after base change by affine schemes mapping into
$Y_j \times_Y X \times_Y X$, we conclude by
Lemma \ref{lemma-factor-through-thickening}.

\medskip\noindent
Let us prove that (5) implies (1). Let $\{Y_j \to Y\}$ be as in (5).
Then we have the fibre product diagram
$$
\xymatrix{
\coprod Y_j \times_Y X \ar[r] \ar[d] &
X \ar[d] \\
\coprod Y_j \times_Y X \times_Y X \ar[r] &
X \times_Y X
}
$$
By assumption the left vertical arrow is quasi-compact or a closed immersion.
It follows from
Spaces, Lemma \ref{spaces-lemma-descent-representable-transformations-property}
that also the right vertical arrow is quasi-compact or a
closed immersion.
\end{proof}

\end{document}
